\documentclass{reportkit}

\setreportkitleftheader{REPORTKIT / ALGORITHM VISUALIZATION TEST}
\setreportkitfooter{Algorithm visualization (interval/heap) acceptance test}
\setreportkitversion{v1.10.0}

\title{ReportKit -- interval and heap algorithm-visualization primitives}
\author{Test build}
\date{17 September 2026}

\begin{document}
\maketitle

\begin{execsummary}
This is an acceptance test for reportkit-algorithm-order.sty's
\texttt{intervalstate} and \texttt{heapstate} primitives (spec sections 11
and 12 of
docs/superpowers/specs/2026-09-16-reportkit-algorithm-visualization-primitives-spec.md).
Both reuse the shared nine-state algorithm vocabulary and the foundational
\texttt{arraystate} cell chrome already exercised in
algorithm\_visuals\_acceptance\_test.tex.
\end{execsummary}

\section{Merge intervals}
\texttt{intervalstate}: three candidate intervals on a shared axis, one
currently overlapping the running merge (\texttt{state=active}), followed by
their merged summary span (\texttt{state=resolved} by default) on its own
row below.

\begin{diagram}[type=interval,width=\textwidth,caption={Merging three overlapping intervals into one summary span.},description={Three intervals on a shared numeric axis: A from one to four, B from three to six overlapping A, and C from eight to ten which does not overlap. A merged summary bar spans one to six below the individual intervals.}]
\begin{intervalstate}
  \interval[state=candidate]{A}{1}{4}
  \interval[state=active]{B}{3}{6}
  \interval[state=candidate]{C}{8}{10}
  \merged{1}{6}
\end{intervalstate}
\end{diagram}

\section{Session gaps and a discarded interval}
A second \texttt{intervalstate} example: a short session below a minimum
duration threshold is marked discarded, distinct from the candidate and
resolved sessions around it.

\begin{diagram}[type=interval,width=\textwidth,caption={Sessionization: one short session is discarded below a minimum-duration threshold.},description={Three sessions on a shared axis: a resolved session from zero to five, a discarded session from six to seven that falls below the minimum duration, and a candidate session from nine to fifteen.}]
\begin{intervalstate}
  \interval[state=resolved]{S1}{0}{5}
  \interval[state=discarded]{S2}{6}{7}
  \interval[state=candidate]{S3}{9}{15}
\end{intervalstate}
\end{diagram}

\section{Heap: backing array plus tree view}
\texttt{heapstate}: a small max-heap declared once via \texttt{\string\values},
with both the backing array and the mirrored binary-tree view drawn
automatically -- tree node positions are derived from array index, never
placed by hand. The root is marked current.

\begin{diagram}[type=heap,width=\textwidth,caption={A seven-element max-heap: backing array and derived tree view.},description={A backing array of seven values -- nine, seven, eight, four, five, six, two -- paired with its binary-tree view. The root, value nine, is marked current.}]
\begin{heapstate}
  \values{9,7,8,4,5,6,2}
  \current{0}
\end{heapstate}
\end{diagram}

\section{Heap: marking an internal node}
\texttt{\string\current} can mark any declared index, not only the root --
here the index-2 node (value eight) is marked current in both its array
cell and its mirrored tree node.

\begin{diagram}[type=heap,width=\textwidth,caption={Marking an internal heap node during a sift-down step.},description={The same seven-element max-heap with the index-two node, value eight, marked current in both the backing array and the tree view.}]
\begin{heapstate}
  \values{9,7,8,4,5,6,2}
  \current{2}
\end{heapstate}
\end{diagram}

\section{Grayscale and semantic check}
\begin{principle}{State is not just color}
\texttt{intervalstate} maps the spec's \texttt{overlap}/\texttt{merged}
concepts onto the shared vocabulary's \texttt{active}/\texttt{resolved}
states rather than inventing new ones, so every interval and heap node still
carries the same border-weight and line-style distinctions as arraystate,
and stays legible if printed in grayscale.
\end{principle}

\end{document}
